\begin{table}[H]
\centering
\footnotesize
\caption{F1 por taxón que explica las variaciones del promedio al cambiar de nivel (validación).}
\label{tab:tax-diagnostics}
\begin{tabularx}{\linewidth}{l l Y r r}
\toprule
\textbf{Modelo} & \textbf{Nivel} & \textbf{Taxón} & \textbf{Obs.} & \textbf{F1}\\
\midrule
R & Clase & Agaricomycetes & 3\,087 & 0,996\\
R & Clase & Pezizomycetes & 11 & 1,000\\
R & Clase & Sordariomycetes & 354 & 0,965\\
R & Filo & Ascomycota & 365 & 0,966\\
R & Filo & Basidiomycota & 3\,087 & 0,996\\
A & Familia & Amanitaceae & 904 & 0,992\\
A & Familia & Auriculariaceae & 288 & 0,962\\
A & Familia & Geastraceae & 79 & 0,949\\
A & Familia & Hericiaceae & 66 & 0,962\\
A & Familia & Lycoperdaceae & 408 & 0,978\\
A & Familia & Morchellaceae & 11 & 0,952\\
A & Familia & Pleurotaceae & 251 & 0,978\\
A & Familia & Polyporaceae & 1\,091 & 0,989\\
A & Familia & Xylariaceae & 354 & 0,986\\
A & Orden & Agaricales & 1\,563 & 0,992\\
A & Orden & Auriculariales & 288 & 0,962\\
A & Orden & Geastrales & 79 & 0,949\\
A & Orden & Pezizales & 11 & 0,952\\
A & Orden & Polyporales & 1\,091 & 0,989\\
A & Orden & Russulales & 66 & 0,962\\
A & Orden & Xylariales & 354 & 0,986\\
\bottomrule
\end{tabularx}
\source{reports/s6/comparison/taxonomic_per_class.csv}
\end{table}
